An axial cross-section is an isosceles triangle with sides $10,10,12$, height $8$, area $48$, and semiperimeter $16$. Its incircle radius, equal to the first sphere radius, is $\dfrac{48}{16}=3$.

If $\theta$ is the half-angle at the cone vertex, $\sin\theta=\dfrac{6}{10}=\dfrac{3}{5}$. A tangent sphere of radius $r$ has center at distance $\dfrac{r}{\sin\theta}$ from the vertex. For adjacent radii $3$ and $r$, external tangency gives
\[
\dfrac{3-r}{\dfrac{3}{5}}=3+r,
\]
so $r=\dfrac{3}{4}$. Thus $\boxed{m+n=7}$.
\begin{asy}
texpreamble("\usepackage{amsmath,amssymb}");
size(240);
defaultpen(fontsize(9pt));
pair O=(0,0),A=(-6,8),B=(6,8);draw(O--A--B--cycle);draw(circle((0,5),3),blue);draw(circle((0,1.25),.75),red);draw(O--(0,8),dashed+gray);label("vertex",O,(0,-1));label("$3$", (1.5,5),(0,1));draw((0,5)--(3,5));label("$r$",(.4,1.25),(0,1));
\end{asy}
